\documentclass[11pt]{article}
\usepackage[a4paper,margin=1in]{geometry}
\usepackage{fontspec}
\usepackage[boldfont=false]{xeCJK}
\setCJKmainfont{FandolSong-Regular.otf}
\usepackage{amsmath}
\usepackage{amssymb}
\usepackage{array}
\usepackage{booktabs}
\usepackage{graphicx}
\usepackage{adjustbox}
\usepackage{enumitem}
\setlist[itemize]{topsep=2pt,partopsep=0pt,parsep=0pt,itemsep=2pt,leftmargin=1.4em}
\setlist[enumerate]{topsep=2pt,partopsep=0pt,parsep=0pt,itemsep=2pt,leftmargin=1.6em}
\usepackage{parskip}
\usepackage{fvextra}
\fvset{breaklines=true,breakanywhere=true,fontsize=\small}
\setlength{\parindent}{0pt}

\begin{document}

\begin{center}
  {\LARGE\bfseries Carbonyl compounds}\\[0.35em]
  {\large A Level Chemistry}
\end{center}
\vspace{0.6em}

\section*{Aldehydes and ketones}

Aldehydes and ketones are \textbf{carbonyl compounds} 羰基化合物 — they contain the C=O \textbf{carbonyl} 羰基 group.

\begin{itemize}
\item in an \textbf{aldehyde} 醛 the carbonyl carbon is on the end of the chain (it also carries an H), written $\text{–CHO}$.

\item in a \textbf{ketone} 酮 the carbonyl carbon is in the middle, between two other carbons.

\end{itemize}

\subsection*{Making aldehydes and ketones}

Both are made by the \textbf{oxidation} 氧化 of an \textbf{alcohol} 醇 with acidified $\text{K}_2\text{Cr}_2\text{O}_7$ or $\text{KMnO}_4$:

\begin{itemize}
\item a \textbf{primary} alcohol, with \textbf{distillation} 蒸馏, gives an aldehyde.

\item a \textbf{secondary} alcohol gives a ketone.

\end{itemize}

\subsection*{Reactions}

\begin{itemize}
\item \textbf{reduction} 还原 with $\text{NaBH}_4$ or $\text{LiAlH}_4$ turns a carbonyl compound back into an alcohol (an aldehyde gives a primary alcohol; a ketone gives a secondary alcohol).

\item reaction with \textbf{hydrogen cyanide} 氰化氢 ($\text{HCN}$), with $\text{KCN}$ as catalyst and heat, adds $\text{H}$ and $\text{CN}$ across the C=O to make a \textbf{hydroxynitrile} 羟基腈. This adds one carbon to the chain.

\end{itemize}

\subsection*{The mechanism: nucleophilic addition}

The reaction with $\text{HCN}$ is a \textbf{nucleophilic addition} 亲核加成. The carbonyl carbon is slightly positive (oxygen pulls the electrons away). So:

\begin{enumerate}
\item the $\text{CN}^-$ ion (a nucleophile) attacks the slightly positive carbon.

\item this breaks the C=O double bond, leaving a negative oxygen ($\text{O}^-$).

\item the $\text{O}^-$ takes an $\text{H}^+$ (from $\text{HCN}$) to finish the hydroxynitrile.

\end{enumerate}

\section*{Tests for carbonyl compounds}

\subsection*{Detecting any carbonyl}

Add \textbf{2,4-DNPH} reagent (2,4-dinitrophenylhydrazine). An orange precipitate confirms that the compound is an aldehyde or a ketone.

\subsection*{Telling an aldehyde from a ketone}

Aldehydes are easily oxidised to carboxylic acids, but ketones are not. Two tests use this difference:

\begin{center}
\adjustbox{max width=\linewidth}{%
\begin{tabular}{lll}
\toprule
\textbf{Test} & \textbf{Aldehyde} & \textbf{Ketone} \\
\midrule
\textbf{Fehling's reagent} 斐林试剂 (blue solution) & turns to a brick-red precipitate & no change \\
\textbf{Tollens' reagent} 托伦试剂 (colourless) & gives a silver mirror & no change \\
\bottomrule
\end{tabular}}
\end{center}

\subsection*{The iodoform test}

If the compound has the $\text{CH}_3\text{CO}-$ group, warming it with alkaline aqueous iodine gives a pale yellow precipitate of \textbf{tri-iodomethane} 三碘甲烷 ($\text{CHI}_3$) and the ion $\text{RCO}_2^-$.


\end{document}
